% ID: SEL-D0028-OB-E01
% ESTADO: cotejada. Segunda pasada visual completa y borrador compilado revisados.
% Fuente: Selectividad/matematicas_aplicadas/suplente junio/Examen-Matemáticas Aplicadas a las Ciencias Sociales.pdf
% SHA-256: a10146fce64a17e21450973c88d66d43571619781f76bc345b5a4b3dee67a16b
% Página: 2; opción B; ejercicio original 1.

Sean las matrices
\[
A=\begin{pmatrix}6&0\\2&4\end{pmatrix},\quad B=\begin{pmatrix}-4\\6\end{pmatrix}\quad\text{y}\quad C=\begin{pmatrix}-2&-2\end{pmatrix}.
\]
\begin{enumerate}[label=\alph*)]
\item \textbf{(1 punto)} Justifique cuáles de las siguientes operaciones se pueden realizar y efectúelas cuando sea posible:
\[
B+2C\cdot A\qquad A-(B\cdot C)^t
\]
\item \textbf{(1.5 puntos)} Resuelva la siguiente ecuación matricial: $\dfrac{1}{5}(B+A\cdot X)=C^t$
\end{enumerate}
